\documentclass[sigconf]{acmart}

\usepackage{cuted}

\AtBeginDocument{%
    \providecommand\BibTeX{{%
        Bib\TeX}}}

\setcopyright{none}
\settopmatter{printacmref=false}
\renewcommand\footnotetextcopyrightpermission[1]{}

\begin{document}

    \title{A Recommender System for Biological Control Agents in Agriculture}

    \author{Daniel Carvalho}
    \email{daniel.mq.carvalho@gmail.com}
    \author{Márcia Barros}
    \authornote{Project mentor}
    \email{marcia.c.a.barros@gmail.com}
    \affiliation{%
        \institution{FCUL}
        \city{Lisbon}
        \country{Portugal}
    }
    \renewcommand{\shortauthors}{Carvalho and Barros}

    \begin{abstract}
        Biological Control Agents (BCAs) offer alternatives to chemical plant protection, but data about their use for
        specific crop and pathogen pairs is spread across regulatory and agricultural sources. Differences in terms
        and data structures hinder BCA identification. We present a pipeline that parses and normalises authorisation
        exports, validates rows with rules, and ranks candidates. The pipeline represents crops, target organisms,
        authorised products, and microbial active substances in a shared schema. Because the sources do not provide
        the same fields, the joined modelling data keeps only rows with complete scientific crop and pathogen pairs.
        The prototype uses these records to separate BCA related rows from other rows and rank candidate BCAs for a
        crop and pathogen query, while keeping products as supporting evidence. The evaluation measures recovery of
        authorisation-derived labels and known regulatory associations, not independent biological efficacy. In a
        warm-start evaluation over 109 crop--pathogen contexts, the hybrid ranker recovered 93.6\% of held-out BCAs
        in the top five inferred candidates, with an MRR of 0.749. These results describe association retrieval from
        the available authorisation data, not biological efficacy.

    \end{abstract}

    \keywords{biological control agents, crop protection, ranking systems, data integration, plant pathogens}

    \maketitle


    \section{Background}\label{sec:background-and-context}

    Pests and plant diseases cause substantial crop losses, so crop protection is important for agricultural
    production. Biological control agents, including beneficial microorganisms and their products, can contribute to
    integrated pest management as alternatives or complements to synthetic pesticides~\cite{glare-biopesticides,
    van-lenteren-biological-control}. Their practical use depends on the crop, target organism, microbial strain,
    formulation, application conditions, and regulatory context. Microbial agents can act through several mechanisms,
    including plant defence induction, competition, hyperparasitism, and antibiosis~\cite{kohl-mode-action}. These
    mechanisms and their dependence on the local environment make evidence from an authorisation record different
    from evidence of efficacy in a new field setting.

    \subsection{Related Work}

    Previous work describes the opportunities and adoption barriers for microbial and other biological control
    agents, including variable field performance, formulation constraints, and the need to integrate them into
    broader pest-management systems~\cite{glare-biopesticides,van-lenteren-biological-control}. Research on
    microbial modes of action also shows that disease suppression can result from interacting mechanisms rather than
    one universal agent--pathogen effect~\cite{kohl-mode-action}. These findings motivate a system that makes
    existing authorisation associations easier to retrieve and inspect while avoiding claims about biological
    performance. The present work therefore focuses on data integration and ranking of traceable candidates.


    \section{Problem and Research Question}\label{sec:motivation-and-project-framing}

    Regulatory sources contain information about BCAs and their authorised uses, but finding an option for a given
    crop and pathogen pair remains difficult. The sources use different structures, terms, and levels of detail.
    They also present regulatory or production records rather than answers to a specific crop protection question.
    We therefore frame BCA identification as a ranking task. The system links crops, pathogens, products, and active
    substances, then ranks candidate BCAs for a query.

    The research question is whether these records can retrieve known BCA associations and rank related candidates for
    a crop--pathogen query. This paper addresses data integration and retrieval. It does not compare efficacy models
    because the current sources contain no independent biological efficacy labels.


    \section{Contributions}\label{sec:research-goals}

    This paper makes \textbf{one contribution}: an authorisation-derived ranking method for BCA candidates in crop--pathogen
    queries when the available data contain no biological efficacy labels. The method represents crop--pathogen
    contexts and BCA values, then keeps product names as evidence that users can inspect.

    We pursued three objectives. First, we integrated heterogeneous authorisation sources into a shared schema while
    preserving unresolved cases for review. Second, we separated authorisation-derived BCA evidence from other rows
    to construct transparent weak labels. Third, we tested whether a hybrid ranker could retrieve held-out BCA
    associations at least as well as direct-support and simple support baselines in warm-start contexts. The
    evaluation hypothesis was that combining direct support with a small co-occurrence signal would match or improve
    retrieval of observed associations.

    The pipeline maps four authorisation sources to a common schema and separates accepted, review, and discarded
    rows. A weakly supervised classifier reproduces these row classes and provides a diagnostic triage score. The
    BCA-level ranker uses accepted Portuguese rows to retrieve known matches and infer candidates from related
    authorisation evidence. We evaluate the classifier on a grouped holdout and the ranker on held-out authorisation
    associations. In the inferred-only view, the hybrid ranker reaches 93.6\% Hit@5 and an MRR of 0.749 over 109
    contexts. These evaluations test label reproduction and association retrieval, not biological efficacy.


    \section{Method}\label{sec:methodology-or-approach}

    The pipeline parses authorisation files into a shared table. It checks the normalised records against microbe and
    pathogen reference data, then separates BCA rows from other rows and from cases that need manual review. These
    classes provide weak labels for an authorisation triage model. The query step ranks candidate BCAs from accepted
    authorisation rows and keeps product records as supporting evidence.

    The system separates three questions:
    \begin{itemize}
        \item \textbf{pipeline:} whether it can interpret a source row consistently;
        \item \textbf{classifier:} whether a complete row resembles accepted BCA authorisation rows or rows requiring review; and
        \item \textbf{query:} which known microbe and product candidates should be shown for a given crop and pathogen.
    \end{itemize}
    None of these steps proves biological effectiveness in a new field context. That requires direct testing. The
    system uses crop-pathogen relationships to identify microbes for further examination.

    \subsection{Pipeline Implementation}\label{subsec:pipeline-implementation}

    To answer this question, we implemented a \textbf{pipeline} that separates source-specific parsing from shared
    processing. Each run produces intermediate and classified records for inspection.

    A registry connects each source to its parser and column mapping. The current sources are the European Union,
    Spain, France, and Portugal. After source-specific mapping, shared classification and output stages handle every
    source consistently.

    Each parser retains original values and maps available fields to canonical crop, pathogen, product, and microbe
    fields. Required columns are checked before processing; a missing required column causes a reported skip rather than
    a partially labelled row. Optional columns remain optional because the exports do not expose the same concepts.

    The source parsers also handle transformations that cannot be expressed as simple renaming. Portuguese rows can
    contain several crop and target names in one cell. The pipeline expands safely paired common and scientific names
    into separate rows. When the two lists have different lengths or cannot be paired, it sends the row to manual
    review. French usage records encode several use attributes in one identifier, which the parser separates into crop,
    application method, and optional target fields. Spanish product, crop, and substance catalogues remain separate
    because the available product data do not provide a registration key that can support a product-to-use join. These
    rules preserve source information while preventing the pipeline from inventing relationships.

    Normalisation removes formatting noise while retaining meaningful qualifiers, and dates keep their date part when
    an export includes a timestamp. Reference and review fields make unresolved values explicit: missing source values
    remain missing, while a failed lookup receives a review reason.

    The output stage separates processing from modelling. Each source receives final BCA, review, discarded, and summary
    records. A second step combines shared columns and keeps only rows with both scientific crop and pathogen names.
    Incomplete rows cannot enter the shared modelling table.
    This design prevents a large product catalogue from being mistaken for a set of verified crop-pathogen uses.

    The ranking and query stages read this data contract rather than source-specific parsers. Tests cover normalisation,
    reference matching, BCA classification, target expansion, combined-data filtering, ranking, query matching, and
    evaluation metrics. They verify intended transformations and output meanings, but do not prove source completeness.

    Source identifiers and output classes remain attached to each row, and expanded Portuguese rows retain the fields
    needed to inspect their original crop and target values. These records support audit and show how parser or
    reference-list changes affect row counts. Processing rules are documented in the accompanying software
    artifact~\cite{project-repository}.

    \subsection{Data Sources and Data Integration}\label{subsec:data-sources-and-data-integration}

    We use four authorisation exports from the European Union, Spain, France, and Portugal. The sources differ in
    language, level of detail, and table structure, so each parser maps its columns to a shared schema. The EU input
    ~\cite{ec-pesticides-database,eu-dataset-snapshot} provides microbial substances, authorisation
    status, approval dates, authorised countries, and functions, but no crop-pathogen authorisation relation. The
    Spanish inputs are product, crop, and substance workbooks~\cite{mapa-regfiweb,spain-dataset-snapshot}; the
    product workbook is a \textbf{product catalogue}, while the crop and substance workbooks remain reference catalogues
    because no registration key links them safely to product uses. The French input is the E-Phy usage export
    ~\cite{anses-ephy,france-dataset-snapshot}. It provides authorised product-use rows and common crop and target
    labels, but not the scientific crop-pathogen pair required by the joined modelling data. The Portuguese input
    consists of several official exports~\cite{dgav-sifito}. These sources provide the most complete crop, pathogen,
    product, authorisation, and microbial substance fields and supply all rows in the current joined modelling data.
    Common fields include crop names, scientific crop names where available, target organisms, scientific pathogen
    names, product names, product functions, authorisation dates, and microbial active substances.

    The Spanish product table does not link products to crops or pests, and the French use export lacks the scientific
    pair fields required by the joined modelling data. The other country records remain available for catalogue
    inspection and audit. Only Portuguese records currently enter the joined modelling data: 1,530 final rows, 10,160
    review rows, and 183,507 discarded rows. The other sources remain outside the ranking catalogue until their missing
    links can be validated.

    Before merging records, the pipeline checks each source for its required columns and skips files that fail the
    check. It keeps optional fields when available. The Portuguese source needs an expansion step because one row can
    list several target names and scientific names. The parser expands safely paired names into crop, pathogen, and
    product records and sends ambiguous pairs to manual review.

    After normalisation, the pipeline adds reference and review fields. It matches microbe and pathogen names against
    curated reference lists after normalising accents, punctuation, and strain wording. EPPO names and codes provide
    an external reference for plant and pest terminology~\cite{eppo-global-database}.
    The pipeline marks a row as BCA related when its microbe value matches the microbe reference data or contains an
    explicit biological indicator such as a strain, virus, bacteriophage, or microorganism. Rows with BCA signal and
    no review reason enter \textit{final}; rows with BCA signal and an unresolved reference or target pair enter
    \textit{review}; all remaining rows enter \textit{discarded}.
    These groups are evidence classes produced by the pipeline, not biological efficacy labels; discarded rows are
    retained for audit and used as negative examples for the classifier.

    \subsection{Ranking Method}\label{subsec:recommendation-methodology}

    We rank the canonical microbial active-substance value. It represents a BCA or a listed BCA mixture. We keep
    products as evidence so users can inspect the authorisation records behind each candidate.

    We build the current ranking catalogue from Portuguese authorisation records only. Portugal is the only
    source that provides complete records for the required fields. We keep the other country outputs separate until we
    recover and validate their missing links.

    We derive the interaction table from accepted \textit{final} rows; \textit{review} rows remain a separate
    unresolved queue. We normalise microbial mixture strings into a canonical order and remove duplicate components.
    Each row represents one crop-pathogen context and one canonical BCA value. We collapse repeated source rows into
    an interaction count and aggregate source, output class, product name, and product function. Support
    uses distinct products when available and interaction counts otherwise. Products remain \textit{evidence}, not ranked items.

    We treat each crop-pathogen context as an implicit user and each BCA as an item. The interaction table becomes a
    context-item matrix. We retain interaction counts for support features but use binary presence for similarity.
    The collaborative component computes cosine similarity between BCA vectors, prefers contexts with the same
    pathogen, and falls back to global contexts when necessary. It shrinks similarities when item overlap is small
    and scores each candidate from its two strongest similarities to the observed query BCAs. If the query has no
    observed items, the collaborative score is zero and direct support provides the fallback.

    The ranker combines context, pathogen, and collaborative support. It also computes crop, total, and dataset
    support, but assigns them zero weight. It normalises count features within the candidate set before weighting them.
    The current weights are $0.475$ for known context
    matches, $0.190$ for context support, $0.285$ for pathogen support, and $0.050$ for collaborative similarity;
    crop, total, and dataset support have weight zero. The query also reports the triage score as a diagnostic, but
    does not use it for ordering in the current evaluation. Its weight is zero and the value remains available for
    audit and manual review.

    The query matches case-insensitive literal text in the scientific crop and common or scientific pathogen fields.
    It separates known matches from inferred candidates. Inferred candidates come from related authorisation evidence;
    they do not establish a new authorisation or a biological efficacy claim.


    \clearpage


    \section{Evaluation Setup}\label{sec:experimental-or-evaluation-setup}

    \begin{strip}
        \centering
        \captionof{table}{Experimental setup and outputs.}
        \label{tab:experimental-setup}
        \scriptsize
        \setlength{\tabcolsep}{3pt}
        \renewcommand{\arraystretch}{1.15}
        \begin{tabular}{p{0.12\textwidth} p{0.20\textwidth} p{0.40\textwidth} p{0.20\textwidth}}
            \hline
            \textbf{Stage} & \textbf{Input}                                                                     & \textbf{Processing or model}                                                                                                          & \textbf{Output}                                                                                                 \\
            \hline
            Pipeline       & Four authorisation exports.                                                        & Source-specific parsing, canonical normalisation, reference matching, and separation into final, review, and discarded rows. & 1,530 final rows, 10,160 review rows, and 183,507 discarded rows in the joined modelling data. \\
            Classifier     & 195,197 joined rows: 11,690 positive final/review rows and 183,507 discarded rows. & TF-IDF text features with balanced logistic regression. An 80/20 grouped holdout is split by dataset, microbe, pathogen, and product. & Accuracy, macro F1, ROC AUC, and average precision of $1.000$. \\
            Ranking        & 1,530 accepted Portuguese final rows.                                              & 595 context-BCA interactions. One edge is held out from each of 109 complete contexts containing at least two BCAs. & Hybrid direct-support and collaborative ranking, with inferred-only results in Table~\ref{tab:ranking-results}. \\
            \hline
        \end{tabular}
    \end{strip}

    We evaluate the pipeline, classifier, and ranker separately. Pipeline checks confirm that the parsers produce the
    expected schema, required columns exist, target expansion produces auditable rows, and each data set produces the
    expected \textit{final}, \textit{review}, and \textit{discarded} files.

    The classifier uses the generated classes as weak labels. Rows in \textit{final} and \textit{review} receive the
    positive label because they contain BCA evidence or need BCA related review. Rows in \textit{discarded} receive the
    negative label because the pipeline found no BCA signal. Only rows with complete scientific crop and pathogen
    pairs enter training. The classifier uses scientific crop and pathogen names, microbe, product, product function,
    and source context; it excludes common crop and pathogen names because their language varies. The TF-IDF
    vectorizer uses unigrams and bigrams, ignores terms seen in only one row, and keeps at most 60,000 features.
    Balanced logistic regression uses the liblinear solver with 1,000 iterations. Reference-match flags,
    extracted genera, review and discard reasons, and purpose text are excluded to reduce direct leakage from the
    pipeline decisions. An 80/20 grouped holdout keeps related dataset, microbe, pathogen, and product rows on the
    same side of the split. The resulting metrics measure reproduction of labels derived from authorisation data.

    For ranking, we select complete crop-pathogen contexts with at least two distinct observed BCAs and hold out one
    deterministic context-BCA edge from each context. We remove the edge before building features, then rank the held
    out BCA among the remaining candidate catalogue. For reproducibility, contexts are ordered by their scientific crop
    and pathogen names and the lexicographically last canonical BCA is selected as the held-out edge. Unavailable targets
    are counted as misses; 107 of 109 inferred-only cases remain covered (98.2\%). We report results for all candidates
    and for an inferred-only view, which removes candidates that still have a known context match and corresponds to the
    query interface.

    We measure ranking quality with HitRate@1, HitRate@5, HitRate@10, mean reciprocal rank (MRR), and normalised
    discounted cumulative gain (NDCG), following standard cumulative-gain evaluation for ranked retrieval~\cite{jarvelin-cumulated-gain}.
    We report candidate coverage separately because the ranker cannot recover a BCA
    that never appears elsewhere in the interaction catalogue in this warm start experiment. The baselines rank BCAs
    by pathogen support and total support.

    This protocol tests retrieval of observed authorisation associations. It does not establish biological efficacy,
    safety, suitability for an untreated crop and pathogen pair, or legal authorisation for a new use.


    \section{Results}\label{sec:results-and-project-outcomes}

    The classifier reproduced the weak labels perfectly: accuracy, macro F1, ROC AUC, and average precision are all
    $1.000$. These values show that it reproduced the pipeline output classes in this split; they do not measure
    biological efficacy.

    Table~\ref{tab:ranking-results} reports the inferred-only results. It separates candidate recovery, measured by
    coverage and HitRate@k, from rank quality, measured by MRR and NDCG@10. Coverage is the fraction of queries in
    which the held-out BCA remains in the candidate catalogue. HitRate@k is the fraction of queries that recover it in
    the first $k$ positions. MRR rewards higher ranks, and NDCG@10 gives less weight to lower-ranked hits.

    \begin{table}[H]
        \caption{Ranking results for the inferred only view ($n=109$ contexts). Higher is better for every metric.}
        \label{tab:ranking-results}
        \centering
        \small
        \setlength{\tabcolsep}{3pt}
        \renewcommand{\arraystretch}{1.15}
        \begin{tabular}{lrrrr}
            \hline
            \multicolumn{5}{l}{\textit{Candidate recovery}} \\
            \textbf{Variant}  & \textbf{Coverage} & \textbf{Hit@1} & \textbf{Hit@5} & \textbf{Hit@10} \\
            \hline
            Hybrid            & 98.2\%            & 60.6\%         & 93.6\%         & 98.2\%          \\
            Direct support    & 98.2\%            & 60.6\%         & 92.7\%         & 98.2\%          \\
            Collaborative     & 98.2\%            & 45.9\%         & 86.2\%         & 93.6\%          \\
            Pathogen baseline & 98.2\%            & 60.6\%         & 92.7\%         & 98.2\%          \\
            Total baseline    & 98.2\%            & 4.6\%          & 47.7\%         & 82.6\%          \\
            \hline
        \end{tabular}
        \vspace{3pt}
        \begin{tabular}{lrr}
            \hline
            \multicolumn{3}{l}{\textit{Rank quality}} \\
            \textbf{Variant}  & \textbf{MRR} & \textbf{NDCG@10} \\
            \hline
            Hybrid            & 0.749        & 0.807            \\
            Direct support    & 0.743        & 0.802            \\
            Collaborative     & 0.645        & 0.715            \\
            Pathogen baseline & 0.743        & 0.802            \\
            Total baseline    & 0.249        & 0.373            \\
            \hline
        \end{tabular}
    \end{table}

    The triage score is not included in the current ranking score. This keeps it available as transparent
    evidence from the authorisation data while avoiding a ranking contribution that was not supported by the held out
    retrieval results. The results should therefore be read as recovery of associations present in the source data,
    not as evidence that one BCA is biologically superior to another.


    \section{Discussion}\label{sec:discussion}

    A pathogen support signal accounts for much of the ranking behaviour in the current data. The sources repeat
    authorisation records around target organisms, while a missing record does not provide a reliable negative label.
    The hybrid ranker slightly improves on direct support alone in HitRate@5, MRR, and NDCG@10. Collaborative support
    alone performs worse. The improvement over direct support is small: $+0.9$ percentage points in Hit@5, $+0.006$
    MRR, and $+0.005$ NDCG@10. No confidence intervals or significance test are available, so these differences should be treated as
    descriptive evidence for a possible small benefit from collaboration rather than as a firm improvement.

    The classifier results also need a narrow interpretation. Grouped holdout keeps related source, microbe, pathogen,
    and product rows on the same side of the split, and rule-derived match and review fields are excluded. However, the
    labels are generated from the same authorisation records and related fields that supply the model features. The
    perfect scores therefore show separability or reproduction of pipeline labels, not generalisation to independently
    judged BCA status or biological outcomes. The discarded class means that no BCA signal was detected; it does not
    mean biological ineffectiveness.

    The ranking evaluation is a warm-start experiment over contexts with several observed BCAs. It does not measure
    performance for a new crop-pathogen pair, a BCA absent from the candidate catalogue, or a future regulatory snapshot.
    The labels inherit the coverage and regulatory bias of the authorisation sources. Because the triage weight is zero,
    the classifier is not part of this ranking evaluation. The experiment uses the authorisation-derived final catalogue
    and does not test end-to-end generalisation.


    \section{Conclusion and Future Work}\label{sec:conclusions-and-future-work}

    This work delivers a pipeline that combines authorisation data from different sources and a BCA ranker for crop
    and pathogen queries. The system combines support counts with context similarity, keeps products as evidence, and
    reports known and inferred candidates separately. The evaluation shows retrieval of observed authorisation
    associations, not biological efficacy. The research question is therefore answered only for warm-start retrieval of
    authorisation-derived associations: the system can recover held-out associations in the current Portuguese catalogue,
    but it cannot establish efficacy, approval, or suitability for new uses.

    Future work should add structured sources, show associated crop and pathogen pairs as evidence, improve BCA mixture
    normalisation, and use a separate validation split to tune ranking weights. It should also test a bounded large
    language model (LLM) stage for irregular datasets. The LLM could suggest column mappings, translate source labels,
    and propose candidate entity matches so that more BCA records can enter the pipeline. Deterministic validation
    and human review must check these suggestions before they reach the modelling table; inferred mappings cannot
    count as verified authorisation evidence. Cold-start and temporal evaluations are also needed before assessing the
    system beyond this authorisation-based prototype.

    \bibliographystyle{../ACM-Reference-Format}
    \bibliography{references}

\end{document}
